\section{Race the risky abstraction against the exact search \hfill \live{}}

\textbf{The claim, and it is not the usual one.} Parallelism here does not make the search faster.
It changes \emph{which algorithms are usable at all}.

Every large technique in this document is sound in one direction only. Environment abstraction
(\S1): if the abstract system is safe, the real one is safe; if the abstract system is unsafe, we
learned nothing and must start over. That one-sidedness is exactly why such techniques are not in
production tools --- a sequential checker has to commit, and a wrong commitment burns the run.

With cores you do not commit. Run the cheap one-sided technique \emph{and} the exact search at the
same time. If the cheap one proves safety first, kill the exact search. If it fails, you spent one
core and the exact search was never interrupted. \textbf{Parallelism converts a risky bet into a
free one.}

\paragraph{We measured an instance of this by accident.} The liveness deadlock gate offers two
strategies that compute the same verdict:

\begin{center}
\begin{tabular}{@{}lrr@{}}
\toprule
model & gate $+$ $a_{LU}$ & exact only \\
\midrule
\texttt{fischer7} & $180\,420$ & $\mathbf{30\,994}$ \\
\texttt{arbiter6} & $\mathbf{28\,591}$ & $66\,393$ \\
\bottomrule
\end{tabular}
\end{center}

$5.8\times$ one way, $2.3\times$ the other, and \emph{which one wins cannot be known in advance} ---
deciding it is as hard as the query. Racing them yields $\min(A,B)$ on every model instead of $A$
always. It is sound because both compute the same verdict, and it is implementable immediately on a
measured spread.

\paragraph{The general version.} Put environment abstraction on the fast lane. Its cost is $O(1)$ in
$N$, so the race costs essentially one idle core, and the upside is the factor of $1000$ from \S1.
This is the specific remedy for what killed the CMC line: the decomposition order was hand-chosen
and a wrong choice wasted the run. Choose wrong in parallel and it costs nothing.

\paragraph{The counterweight, stated because it bounds the whole idea.} Parallelism also actively
\emph{hurts} state count when it is used the wrong way. Behrmann et al.\ measured \textsc{Dacapo}
rising from $45\,001$ states on one node to $110\,634$ on seventeen, purely from parallelising a
subsumption search: a state covered in one order is expanded in another. We hit the same wall from
the other side --- \verifyta{} \texttt{-o\{0,1,2\}} gives identical counts, DFS is $31.6\%$ worse on
\texttt{phil}, and \texttt{fischer8} is bit-identical from $1$ to $12$ threads \emph{only because}
strict level synchronisation is enforced.

So the rule this document ends on:

\begin{center}
\emph{Use cores for redundant, independently sound work. Never use them to relax the search order.}
\end{center}

One direction is free. The other is a tax.
